\chapter{数学预备知识}
\label{chap:math-foundations}

本章准备后文反复使用的工具：复数、函数相等的几种含义、内积与正交投影。先掌握用法，细节可回头再补。

\section{复数、复指数与欧拉公式}
\label{sec:complex-numbers}

复数用一个数同时表示振幅与相位，频域分析中会反复使用。虚数单位记作 $\iu$：
\[
  \iu^2=-1.
\]
工程中采用 $\mathrm{j}$，以免与电流符号 $\mathrm{i}$ 混淆。

\subsection{代数形式与复平面}

\begin{definition}[复数的代数形式]
  \keyterm{复数} $z$ 写成
  \[
    z=x+\iu y, \qquad x,y\in\R,
  \]
  $x=\operatorname{Re}\{z\}$ 是实部，$y=\operatorname{Im}\{z\}$ 是虚部。两个复数相等，当且仅当\red{实部、虚部分别相等}。
\end{definition}

$z=x+\iu y$ 对应复平面上的点 $(x,y)$：横轴为实轴，纵轴为虚轴。一个复数由此分成两部分：到原点的距离和所指方向。

\begin{definition}[模与幅角]
  复数 $z=x+\iu y$ 的\keyterm{模（magnitude）}为
  \begin{equation}
    |z|=\sqrt{x^2+y^2}.
    \label{eq:complex-magnitude}
  \end{equation}
  当 $z\ne 0$ 时，满足
  \[
    x=|z|\cos\theta,\qquad y=|z|\sin\theta
  \]
  的角 $\theta$ 叫 $z$ 的\keyterm{幅角（argument）}。幅角不唯一：
  \[
    \arg z=\theta+2k\pi,\qquad k\in\Z.
  \]
  区间 $(-\pi,\pi]$ 内的那一个叫幅角主值 $\operatorname{Arg}z$。零的模为零，\red{幅角没有定义}。
\end{definition}

由 $x,y$ 求角时，\hl{先判断象限}，再取反正切。程序中直接使用双参数函数 $\operatorname{atan2}(y,x)$。

\begin{pitfall}[只用反正切求幅角]
  $\arctan(y/x)$ 无法区分相差 $\pi$ 的两个方向，$x=0$ 时又无定义。$1+\iu$ 与 $-1-\iu$ 的虚实部之比相同，主幅角却分别是 $\pi/4$ 和 $-3\pi/4$。\hl{先判断象限}，程序中用 $\operatorname{atan2}(y,x)$。
\end{pitfall}

\subsection{共轭与基本运算}

\begin{definition}[复共轭]
  $z=x+\iu y$ 的\keyterm{复共轭（complex conjugate）}为
  \[
    z^*=x-\iu y.
  \]
  几何上是\red{关于实轴镜像}。$z^*$ 也写作 $\bar z$。
\end{definition}

常用关系：
\begin{align*}
  z+z^*&=2\operatorname{Re}\{z\},
  &z-z^*&=2\iu\operatorname{Im}\{z\},\\
  zz^*&=|z|^2,
  &\frac{1}{z}&=\frac{z^*}{|z|^2}\quad(z\ne0),\\
  (z_1z_2)^*&=z_1^*z_2^*,
  &|z_1z_2|&=|z_1|\,|z_2|.
\end{align*}
其中 $|z_1z_2|=|z_1|\,|z_2|$ 说明复数相乘时模相乘。

\subsection{欧拉公式与极坐标形式}

\keyterm{欧拉公式（Euler's formula）}把复指数与三角函数联系起来：
\begin{equation}
  \e^{\iu\theta}=\cos\theta+\iu\sin\theta.
  \label{eq:euler-formula}
\end{equation}
任意非零复数都可写成
\begin{equation}
  z=r\e^{\iu\theta},
  \qquad r=|z|>0,
  \qquad \theta\in\arg z.
  \label{eq:complex-polar-form}
\end{equation}
这称为\keyterm{极坐标形式}。$r$ 是长度，$\e^{\iu\theta}$ 是单位圆上方向为 $\theta$ 的点。

乘以 $\e^{\iu\theta}$ 相当于在复平面上\red{逆时针旋转 $\theta$}，长度不变；再乘以正实数 $r$，长度变为原来的 $r$ 倍。复数乘法因此是\red{先缩放再旋转}。

若 $z_1=r_1\e^{\iu\theta_1}$、$z_2=r_2\e^{\iu\theta_2}$，则
\begin{align*}
  z_1z_2&=r_1r_2\e^{\iu(\theta_1+\theta_2)},\\
  \frac{z_1}{z_2}&=\frac{r_1}{r_2}\e^{\iu(\theta_1-\theta_2)}
  \qquad(z_2\ne0).
\end{align*}
相乘：\red{模相乘、幅角相加}；相除：\red{模相除、幅角相减}。角度都按模 $2\pi$ 理解。

整数 $n$ 的\keyterm{棣莫弗公式}：
\[
  \bigl(r\e^{\iu\theta}\bigr)^n
  =r^n\e^{\iu n\theta}
  =r^n\bigl(\cos n\theta+\iu\sin n\theta\bigr).
\]
反过来，非零复数 $z=r\e^{\iu\theta}$ 的 $n$ 次方根有 $n$ 个：
\begin{equation}
  z_k=r^{1/n}\e^{\iu(\theta+2k\pi)/n},
  \qquad k=0,1,\ldots,n-1.
  \label{eq:complex-nth-roots}
\end{equation}
它们模相同，角度依次差 $2\pi/n$，\red{均匀分布在同一个圆上}。后文讨论零极点分布时会遇到这一结构。

\begin{example}[代数形式转指数形式]
  $z=-1+\iu\sqrt{3}$。由\cref{eq:complex-magnitude}得 $|z|=2$。点 $(-1,\sqrt{3})$ 在第二象限，所以
  \[
    \operatorname{Arg}z=\frac{2\pi}{3},
    \qquad
    z=2\e^{\iu 2\pi/3}.
  \]
  也可写成 $2\e^{\iu(2\pi/3+2k\pi)}$，$k\in\Z$。
\end{example}

在欧拉公式中分别取 $\theta$ 与 $-\theta$，相加、相减，得到信号分析中最常用的两组关系：
\begin{align}
  \cos\theta
    &=\frac{\e^{\iu\theta}+\e^{-\iu\theta}}{2},
  \label{eq:cos-complex-exponential}\\
  \sin\theta
    &=\frac{\e^{\iu\theta}-\e^{-\iu\theta}}{2\iu}.
  \label{eq:sin-complex-exponential}
\end{align}
实正弦、余弦都可分解为\red{一对反向旋转的复指数}。

\subsection{复指数与正弦信号}

设 $s=\sigma+\iu\omega$，则
\begin{equation}
  \e^{st}=\e^{\sigma t}\e^{\iu\omega t}
  =\e^{\sigma t}\bigl(\cos\omega t+\iu\sin\omega t\bigr).
  \label{eq:complex-exponential-signal}
\end{equation}
$\e^{\sigma t}$ 决定包络，$\omega t$ 决定相位：$\sigma<0$ 时\red{幅度衰减}，$\sigma>0$ 时\red{幅度增长}，$\sigma=0$ 时在复平面上匀速旋转。$\omega\ne0$ 时，$\e^{\iu\omega t}$ 的基本周期为
\[
  T_0=\frac{2\pi}{|\omega|}.
\]

带初相位的实正弦也能写成复指数的实部。$A\geq0$ 时，
\begin{equation}
  A\cos(\omega t+\varphi)
  =\operatorname{Re}\!\left\{A\e^{\iu\varphi}\e^{\iu\omega t}\right\}.
  \label{eq:sinusoid-complex-form}
\end{equation}
复数 $A\e^{\iu\varphi}$ 同时保存振幅与初相位。计算正弦稳态响应时，\hl{先算右边的复指数，最后取实部}。

离散时间复指数：
\[
  x[n]=z^n=r^n\e^{\iu\omega n},\qquad z\ne0,\quad n\in\Z.
\]
$r$ 决定增长或衰减，$\omega$ 决定每前进一步转过的角度。$n$ 只取整数，所以 $\omega$ 与 $\omega+2k\pi$ \red{给出完全相同的序列}。$\e^{\iu\omega n}$ 是周期序列，当且仅当存在整数 $N,m$ 使
\[
  \omega N=2\pi m,
\]
也就是 $\omega/(2\pi)$ 为有理数。最小的正整数 $N$ 是基本周期。

\begin{pitfall}[把连续时间的周期公式直接搬过来]
  连续复指数只要 $\omega\ne0$ 就有周期；离散还要求 $\omega/(2\pi)$ 是有理数。$\e^{\iu n}$ 不满足这一条件，\red{没有整数基本周期}。
\end{pitfall}

相位只在模 $2\pi$ 意义下确定。$\varphi$ 与 $\varphi+2k\pi$ 是同一个方向。比较相位时\blue{先说清主值区间}；数值在 $\pi$ 和 $-\pi$ 附近发生跳变，并不表示信号突然改变。

相位变化 $\varphi$ 可换算成时间平移 $t_0=-\varphi/\omega$：
\[
  \cos(\omega t+\varphi)=\cos\bigl(\omega(t-t_0)\bigr).
\]
同一个相位差，\red{角频率不同，对应的时间平移就不同}。“相差 $\varphi$”并不等于“延迟 $\varphi$”。

\section{函数相等、几乎处处相等与广义函数}
\label{sec:function-equality-and-distributions}

“两个函数相等”有几种不同含义。若不先说明采用哪一种，阶跃在跳变点的取值、冲激的乘法就容易混淆。

\subsection{逐点相等与几乎处处相等}

\begin{definition}[逐点相等]
  对定义域内每个 $t$ 都有 $f_1(t)=f_2(t)$，叫\keyterm{逐点相等（pointwise equal）}。\red{只要有一个点不同}，就不是逐点相等。
\end{definition}

\begin{definition}[几乎处处相等]
  若 $f_1$ 与 $f_2$ 只在一个零测集上不同，则称\keyterm{几乎处处相等（equal almost everywhere）}，记
  \[
    f_1=f_2\quad\text{a.e.}
  \]
  有限集、可数集都是零测集，所以\red{有限个点的差异不影响几乎处处相等}。
\end{definition}

单位阶跃是现成的例子。有的教材把 $\stepfun(0)$ 定为 $0$，有的定为 $1/2$ 或 $1$。它们不逐点相等，但普通区间上的积分结果相同。\blue{只涉及普通积分时}，\red{单个端点值可以忽略}。

\begin{pitfall}[把几乎处处相等当成处处相等]
  几乎处处相等在积分意义下没有差别，但不表示任意点都可直接代入。点值算子只读取 $f(0)$，冲激只抽取所在点的函数值。遇到这类运算，\blue{先判断它是否依赖单点取值}。
\end{pitfall}

\subsection{广义函数按积分作用相等}

冲激 $\impulse(t)$ 是\keyterm{广义函数（generalized function）}，严格地说属于分布。它没有可逐点比较的普通函数值，判等只能依据积分作用：若两个广义函数 $F_1,F_2$ 对任意\blue{足够光滑、在有限范围外为零}的测试函数 $\varphi$ 都满足
\begin{equation}
  \int_{-\infty}^{\infty}\varphi(t)F_1(t)\dd t
  =\int_{-\infty}^{\infty}\varphi(t)F_2(t)\dd t,
  \label{eq:distribution-equality}
\end{equation}
就认为 $F_1=F_2$。\red{不比较函数值，只比较对测试函数的积分结果}。

本课程只使用这一条，不展开测度论。后文的冲激恒等式均通过对测试函数积分验证。

\section{内积与正交展开基础}
\label{sec:inner-product-foundations}

傅里叶级数、傅里叶变换本质上都是正交展开。先说明内积与投影，后文可直接套用公式。

复信号空间里，区间 $[a,b]$ 上的\keyterm{内积}定义为
\begin{equation}
  \langle x,y\rangle
  =\int_a^b x(t)\overline{y(t)}\dd t,
  \label{eq:continuous-signal-inner-product}
\end{equation}
离散序列则为
\begin{equation}
  \langle x,y\rangle
  =\sum_{k=m}^{n}x[k]\overline{y[k]}.
  \label{eq:discrete-signal-inner-product}
\end{equation}
\red{共轭放在第二个因子上}，于是有共轭对称、对第一个因子线性、正定：
\begin{align*}
  \langle x,y\rangle&=\overline{\langle y,x\rangle},&
  \langle \lambda x,y\rangle&=\lambda\langle x,y\rangle,\\
  \langle x+z,y\rangle&=\langle x,y\rangle+\langle z,y\rangle,&
  \langle x,x\rangle&=\int_a^b|x(t)|^2\dd t\geq0.
\end{align*}
$\langle x,y\rangle=0$ 时称 $x,y$ 正交。复信号中\red{共轭不可省略}，否则 $\langle x,x\rangle$ 可能不是非负实数。

两两正交、又能表示所考虑空间中任意信号的函数族 $\{e_k(t)\}$，叫\keyterm{正交基}。把它写成
\begin{equation}
  x(t)=\sum_k a_k e_k(t).
  \label{eq:orthogonal-expansion}
\end{equation}
两边与 $e_r$ 取内积。正交性把 $k\ne r$ 的项全消掉：
\[
  \langle x,e_r\rangle
  =\sum_k a_k\langle e_k,e_r\rangle
  =a_r\langle e_r,e_r\rangle.
\]
于是
\begin{equation}
  \boxed{a_r=\frac{\langle x,e_r\rangle}{\langle e_r,e_r\rangle}}.
  \label{eq:orthogonal-projection-coefficient}
\end{equation}
正交归一基满足 $\langle e_r,e_r\rangle=1$，系数为 $a_r=\langle x,e_r\rangle$。\red{展开系数即信号沿各基方向的投影}。

给定一组线性无关但不正交的 $\{\alpha_k\}$，\keyterm{Gram--Schmidt 方法}\hl{逐个减去已有方向上的投影}：
\begin{align*}
  \beta_1&=\alpha_1,\\
  \beta_2&=\alpha_2-
    \frac{\langle\alpha_2,\beta_1\rangle}{\langle\beta_1,\beta_1\rangle}\beta_1,\\
  \beta_3&=\alpha_3-
    \frac{\langle\alpha_3,\beta_1\rangle}{\langle\beta_1,\beta_1\rangle}\beta_1-
    \frac{\langle\alpha_3,\beta_2\rangle}{\langle\beta_2,\beta_2\rangle}\beta_2.
\end{align*}
计算 $\beta_3$ 时，\hl{第二个投影方向应使用已算出的} $\beta_2$。一般递推为
\begin{align}
  \beta_1&=\alpha_1,\\
  \beta_k&=\alpha_k-
  \sum_{m=1}^{k-1}
  \frac{\langle\alpha_k,\beta_m\rangle}
       {\langle\beta_m,\beta_m\rangle}\beta_m,
  \qquad k\geq2.
  \label{eq:gram-schmidt}
\end{align}
最后令 $e_k=\beta_k/\sqrt{\langle\beta_k,\beta_k\rangle}$，即得正交归一基。后文只需将具体的傅里叶函数\red{代入这套投影框架}。

% 后续按资料进度补充“线性代数与正交展开视角”以及
% “微分方程、差分方程与收敛基础”；未成文前不生成空节标题。
